% !TeX root = ../main.tex
\section{Conclusion}
\label{sec:conclusion}

Guitar tablature transcription has moved from frame-wise CNNs on acoustic
guitar to modular systems with learned fingering and technique stages, and to
generative end-to-end models. In parallel, large sequence-to-sequence
transcribers trained on real multi-instrument data now handle full mixes and
distorted electric guitar, but stop at notes. For the use case of
transcribing a guitar part from a commercial song, the practical state of the
art is a modular pipeline, and our benchmark supports it: a band-split
separator brings note F1 on mixes from 0.50 to 0.81, a guitar model trained
on many tones or MuScriptor handles amplified guitar, and the resulting
tablature F1 on electric guitar (0.55--0.63) is well above end-to-end TabCNN
(0.37--0.40). Its weakest stage is fingering on melodic lines, where learned
models trained on tablature (MIDI-to-Tab, Fretting-Transformer) or conditioned
on audio (TART) have not yet been released for comparison. The most promising
direction is to close the split between the two lines of work, by giving
mix-level transcribers string and fret outputs, and to build a benchmark of
real recordings with reference tablatures so that progress on the actual task
can be measured.
